\problem{50}{(สอวน ม.4 2552) รถไฟขบวนหนึ่งซึ่งปกติแล้วแล่นด้วยอัตราเร็ว $60\ \mathrm{km/h}$ ตลอดระยะทางหนึ่งของการเดินทาง พบว่ามีการซ่อมทางเกิดขึ้น ทำให้ต้องลดอัตราเร็วเป็น $15\ \mathrm{km/h}$ ในช่วงระยะทาง $1.0\ \mathrm{km}$ แล้วเคลื่อนที่ด้วยอัตราเร็วคงตัวที่ลดลงนี้เป็นระยะทาง $1.0\ \mathrm{km}$ ตลอดทางที่กำลังซ่อมอยู่ และเร่งเครื่องให้กลับมีอัตราเร็วเท่าเดิมในระยะทางอีก $1.5\ \mathrm{km}$ ให้สมมติว่าความเร่งของรถไฟทั้งตอนที่เคลื่อนที่ช้าลงและตอนที่เคลื่อนที่เร็วขึ้นนั้นมีขนาดคงตัว (แต่ไม่จำเป็นต้องมีขนาดเท่ากัน) จงหาว่ารถไฟเสียเวลาไปมากกว่าตอนที่แล่นปกตินานเท่าใด}
